\documentclass[10pt,a4paper]{amsart}
\usepackage{amsmath,amssymb,amsthm,graphicx}
% For commutative diagrams we use Xy-pic
\usepackage[all]{xy}

% We use 2cell for 2-commutative diagrams.
\xyoption{2cell}
% Theorem environments.
%
\theoremstyle{plain}
\newtheorem{theorem}[subsection]{Theorem}
\newtheorem{proposition}[subsection]{Proposition}
\newtheorem{lemma}[subsection]{Lemma}

\theoremstyle{definition}
\newtheorem{definition}[subsection]{Definition}
\newtheorem{example}[subsection]{Example}
\newtheorem{exercise}[subsection]{Exercise}
\newtheorem{situation}[subsection]{Situation}

\theoremstyle{remark}
\newtheorem{remark}[subsection]{Remark}
\newtheorem{remarks}[subsection]{Remarks}

\numberwithin{equation}{subsection}

% Macros
%
\def\lim{\mathop{\mathrm{lim}}\nolimits}
\def\colim{\mathop{\mathrm{colim}}\nolimits}
\def\Spec{\mathop{\mathrm{Spec}}}
\def\Hom{\mathop{\mathrm{Hom}}\nolimits}
\def\Ext{\mathop{\mathrm{Ext}}\nolimits}
\def\SheafHom{\mathop{\mathcal{H}\!\mathit{om}}\nolimits}
\def\SheafExt{\mathop{\mathcal{E}\!\mathit{xt}}\nolimits}
\def\Sch{\mathit{Sch}}
\def\Mor{\mathop{\mathrm{Mor}}\nolimits}
\def\Ob{\mathop{\mathrm{Ob}}\nolimits}
\def\Sh{\mathop{\mathit{Sh}}\nolimits}
\def\NL{\mathop{N\!L}\nolimits}
\def\CH{\mathop{\mathrm{CH}}\nolimits}
\def\proetale{{pro\text{-}\acute{e}tale}}
\def\etale{{\acute{e}tale}}
\def\QCoh{\mathit{QCoh}}
\def\Ker{\mathop{\mathrm{Ker}}}
\def\Im{\mathop{\mathrm{Im}}}
\def\Coker{\mathop{\mathrm{Coker}}}
\def\Coim{\mathop{\mathrm{Coim}}}

% Boxtimes
%
\DeclareMathSymbol{\boxtimes}{\mathbin}{AMSa}{"02}

%
% Macros for moduli stacks/spaces
%
\def\QCohstack{\mathcal{QC}\!\mathit{oh}}
\def\Cohstack{\mathcal{C}\!\mathit{oh}}
\def\Spacesstack{\mathcal{S}\!\mathit{paces}}
\def\Quotfunctor{\mathrm{Quot}}
\def\Hilbfunctor{\mathrm{Hilb}}
\def\Curvesstack{\mathcal{C}\!\mathit{urves}}
\def\Polarizedstack{\mathcal{P}\!\mathit{olarized}}
\def\Complexesstack{\mathcal{C}\!\mathit{omplexes}}
% \Pic is the operator that assigns to X its picard group, usage \Pic(X)
% \Picardstack_{X/B} denotes the Picard stack of X over B
% \Picardfunctor_{X/B} denotes the Picard functor of X over B
\def\Pic{\mathop{\mathrm{Pic}}\nolimits}
\def\Picardstack{\mathcal{P}\!\mathit{ic}}
\def\Picardfunctor{\mathrm{Pic}}
\def\Deformationcategory{\mathcal{D}\!\mathit{ef}}
\newenvironment{reference}{}{}
\newenvironment{slogan}{}{}
\newenvironment{history}{}{}
\usepackage{cleveref}
\crefname{originallemma}{original lemma}{original lemmas}
\Crefname{originallemma}{Original lemma}{Original lemmas}
\makeatletter
\newlabel{more-algebra-lemma-homotopy-koszul}{{0663}{source}{}{sourcepointer-1}{}}
\newlabel{more-algebra-lemma-homotopy-koszul@cref}{{[originallemma][0][]0663}{[1][0][]source}{}{}{}}
\makeatother
\crefname{originalsection}{original section}{original sections}
\Crefname{originalsection}{Original section}{Original sections}
\makeatletter
\newlabel{homology-section-filtered-complex}{{012K}{source}{}{sourcepointer-2}{}}
\newlabel{homology-section-filtered-complex@cref}{{[originalsection][0][]012K}{[1][0][]source}{}{}{}}
\makeatother
\crefname{originallemma}{original lemma}{original lemmas}
\Crefname{originallemma}{Original lemma}{Original lemmas}
\makeatletter
\newlabel{homology-lemma-filtered-complex-ss-converges}{{012V}{source}{}{sourcepointer-3}{}}
\newlabel{homology-lemma-filtered-complex-ss-converges@cref}{{[originallemma][0][]012V}{[1][0][]source}{}{}{}}
\makeatother
\crefname{originallemma}{original lemma}{original lemmas}
\Crefname{originallemma}{Original lemma}{Original lemmas}
\makeatletter
\newlabel{algebra-lemma-map-AR}{{00IO}{source}{}{sourcepointer-4}{}}
\newlabel{algebra-lemma-map-AR@cref}{{[originallemma][0][]00IO}{[1][0][]source}{}{}{}}
\makeatother
\crefname{originalsection}{original section}{original sections}
\Crefname{originalsection}{Original section}{Original sections}
\makeatletter
\newlabel{algebra-section-Noetherian-local}{{00K4}{source}{}{sourcepointer-5}{}}
\newlabel{algebra-section-Noetherian-local@cref}{{[originalsection][0][]00K4}{[1][0][]source}{}{}{}}
\makeatother
\crefname{originallemma}{original lemma}{original lemmas}
\Crefname{originallemma}{Original lemma}{Original lemmas}
\makeatletter
\newlabel{lemma-leading-coefficient}{{0AZY}{source}{}{sourcepointer-6}{}}
\newlabel{lemma-leading-coefficient@cref}{{[originallemma][0][]0AZY}{[1][0][]source}{}{}{}}
\makeatother
\crefname{originaldefinition}{original definition}{original definitions}
\Crefname{originaldefinition}{Original definition}{Original definitions}
\makeatletter
\newlabel{algebra-definition-ideal-definition}{{07DU}{source}{}{sourcepointer-7}{}}
\newlabel{algebra-definition-ideal-definition@cref}{{[originaldefinition][0][]07DU}{[1][0][]source}{}{}{}}
\makeatother
\crefname{originalproposition}{original proposition}{original propositions}
\Crefname{originalproposition}{Original proposition}{Original propositions}
\makeatletter
\newlabel{algebra-proposition-hilbert-function-polynomial}{{00K8}{source}{}{sourcepointer-8}{}}
\newlabel{algebra-proposition-hilbert-function-polynomial@cref}{{[originalproposition][0][]00K8}{[1][0][]source}{}{}{}}
\makeatother
\crefname{originallemma}{original lemma}{original lemmas}
\Crefname{originallemma}{Original lemma}{Original lemmas}
\makeatletter
\newlabel{algebra-lemma-support-dimension-d}{{00L8}{source}{}{sourcepointer-9}{}}
\newlabel{algebra-lemma-support-dimension-d@cref}{{[originallemma][0][]00L8}{[1][0][]source}{}{}{}}
\makeatother
\title{Hilbert–Samuel multiplicity equals the Koszul Euler characteristic}
\author{The Stacks Project Authors}
\date{stacks-project:intersection.tex:local-snapshot-20261011; GNU FDL 1.2 or later}
\begin{document}
\maketitle
\noindent\textit{Editorial note.} The historical review incorrectly paraphrases equality of Euler sums as E\_infinity=E\_1. The source 1064–1069 only equates alternating length sums; no degeneration at E\_1 is assumed.\par\medskip
\noindent\textit{Editorial note.} The finite-difference identity is ordinary polynomial algebra and is an allowed prerequisite; it is not another counted problem.\par\medskip
\noindent\textit{Allowed standard prerequisites (editorial).}\begin{itemize}
\item Explicit prerequisite boundary. Noetherian local algebra, Hilbert–Samuel polynomial, Artin–Rees, filtered-complex spectral sequences, length additivity, and the elementary finite-difference identity. None states the Koszul multiplicity equality.
\end{itemize}\medskip
\ifdefined\hypertarget\hypertarget{sourcepointer-1}{}\fi
\noindent\textit{Source reference.} Original lemma 0663: focused-seeds/AGP/sources/stacks-snapshot-20261011/tags.txt, source lines 7922--7922. Stable Stacks tag for an expressly allowed foundational or notational input; the counted construction remains the two concrete core obligations. The retained local core is not replaced by this pointer. See standard\_prerequisites.\par\medskip
\ifdefined\hypertarget\hypertarget{sourcepointer-2}{}\fi
\noindent\textit{Source reference.} Original section 012K: focused-seeds/AGP/sources/stacks-snapshot-20261011/tags.txt, source lines 1402--1402. Stable Stacks tag for an expressly allowed foundational or notational input; the counted construction remains the two concrete core obligations. The retained local core is not replaced by this pointer. See standard\_prerequisites.\par\medskip
\ifdefined\hypertarget\hypertarget{sourcepointer-3}{}\fi
\noindent\textit{Source reference.} Original lemma 012V: focused-seeds/AGP/sources/stacks-snapshot-20261011/tags.txt, source lines 1413--1413. Stable Stacks tag for an expressly allowed foundational or notational input; the counted construction remains the two concrete core obligations. The retained local core is not replaced by this pointer. See standard\_prerequisites.\par\medskip
\ifdefined\hypertarget\hypertarget{sourcepointer-4}{}\fi
\noindent\textit{Source reference.} Original lemma 00IO: focused-seeds/AGP/sources/stacks-snapshot-20261011/tags.txt, source lines 686--686. Stable Stacks tag for an expressly allowed foundational or notational input; the counted construction remains the two concrete core obligations. The retained local core is not replaced by this pointer. See standard\_prerequisites.\par\medskip
\ifdefined\hypertarget\hypertarget{sourcepointer-5}{}\fi
\noindent\textit{Source reference.} Original section 00K4: focused-seeds/AGP/sources/stacks-snapshot-20261011/tags.txt, source lines 738--738. Stable Stacks tag for an expressly allowed foundational or notational input; the counted construction remains the two concrete core obligations. The retained local core is not replaced by this pointer. See standard\_prerequisites.\par\medskip
\ifdefined\hypertarget\hypertarget{sourcepointer-6}{}\fi
\noindent\textit{Source reference.} Original lemma 0AZY: focused-seeds/AGP/sources/stacks-snapshot-20261011/tags.txt, source lines 13841--13841. Stable Stacks tag for an expressly allowed foundational or notational input; the counted construction remains the two concrete core obligations. The retained local core is not replaced by this pointer. See standard\_prerequisites.\par\medskip
\ifdefined\hypertarget\hypertarget{sourcepointer-7}{}\fi
\noindent\textit{Source reference.} Original definition 07DU: focused-seeds/AGP/sources/stacks-snapshot-20261011/tags.txt, source lines 9424--9424. Stable Stacks tag for an expressly allowed foundational or notational input; the counted construction remains the two concrete core obligations. The retained local core is not replaced by this pointer. See standard\_prerequisites.\par\medskip
\ifdefined\hypertarget\hypertarget{sourcepointer-8}{}\fi
\noindent\textit{Source reference.} Original proposition 00K8: focused-seeds/AGP/sources/stacks-snapshot-20261011/tags.txt, source lines 742--742. Stable Stacks tag for an expressly allowed foundational or notational input; the counted construction remains the two concrete core obligations. The retained local core is not replaced by this pointer. See standard\_prerequisites.\par\medskip
\ifdefined\hypertarget\hypertarget{sourcepointer-9}{}\fi
\noindent\textit{Source reference.} Original lemma 00L8: focused-seeds/AGP/sources/stacks-snapshot-20261011/tags.txt, source lines 778--778. Stable Stacks tag for an expressly allowed foundational or notational input; the counted construction remains the two concrete core obligations. The retained local core is not replaced by this pointer. See standard\_prerequisites.\par\medskip
\begin{theorem}
\label{theorem-multiplicity-with-koszul}
\begin{reference}
\cite[Theorem 1 in part B of Chapter IV]{Serre_algebre_locale}
\end{reference}
Let $A$ be a Noetherian local ring. Let $I = (f_1, \ldots, f_r) \subset A$
be an ideal of definition. Let $M$ be a finite $A$-module. Then
$$
e_I(M, r) = \sum
(-1)^i\text{length}_A H_i(K_\bullet(f_1, \ldots, f_r) \otimes_A M)
$$
\end{theorem}



\section{Algebraic multiplicities}
\label{section-multiplicities}

\noindent
Let $(A, \mathfrak m, \kappa)$ be a Noetherian local ring.
Let $M$ be a finite $A$-module and let $I \subset A$ be an ideal
of definition (Algebra, Definition \ref{algebra-definition-ideal-definition}).
Recall that the function
$$
\chi_{I, M}(n) = \text{length}_A(M/I^nM) =
\sum\nolimits_{p = 0, \ldots, n - 1} \text{length}_A(I^pM/I^{p + 1}M)
$$
is a numerical polynomial
(Algebra, Proposition \ref{algebra-proposition-hilbert-function-polynomial}).
The degree of this polynomial is equal to $\dim(\text{Supp}(M))$ by
Algebra, Lemma \ref{algebra-lemma-support-dimension-d}.

\begin{definition}
\label{definition-multiplicity}
In the situation above, assume $d \geq \dim(\text{Supp}(M))$.
In this case, if  $d > \dim(\text{Supp}(M))$, then we set $e_I(M, d) = 0$
and if $d = \dim(\text{Supp}(M))$, then we set $e_I(M, d)$ equal to $d!$
times the leading coefficient of the numerical polynomial $\chi_{I, M}$.
Thus in both cases we have
$$
\chi_{I, M}(n) \sim e_I(M, d) \frac{n^d}{d!} + \text{lower order terms}
$$
The {\it multiplicity of $M$ for the ideal of definition $I$}
is $e_I(M) = e_I(M, \dim(\text{Supp}(M)))$.
\end{definition}

\begin{proof}
Let us change the Koszul complex $K_\bullet(f_1, \ldots, f_r)$ into a cochain
complex $K^\bullet$ by setting $K^n = K_{-n}(f_1, \ldots, f_r)$.
Then $K^\bullet$ is sitting in degrees $-r, \ldots, 0$ and
$H^i(K^\bullet \otimes_A M) = H_{-i}(K_\bullet(f_1, \ldots, f_r) \otimes_A M)$.
The statement of the theorem makes sense as the modules
$H^i(K^\bullet \otimes M)$ are annihilated by $f_1, \ldots, f_r$
(More on Algebra, Lemma \ref{more-algebra-lemma-homotopy-koszul})
hence have finite length.
Define a filtration on the complex $K^\bullet$ by setting
$$
F^p(K^n \otimes_A M) =
I^{\max(0, p + n)}(K^n \otimes_A M),\quad p \in \mathbf{Z}
$$
Since $f_i I^p \subset I^{p + 1}$ this is a filtration by subcomplexes.
Thus we have a filtered complex and we obtain a spectral sequence, see
Homology, Section \ref{homology-section-filtered-complex}.
We have
$$
E_0 = \bigoplus\nolimits_{p, q} E_0^{p, q} =
\bigoplus\nolimits_{p, q} \text{gr}^p(K^{p + q} \otimes_A M) =
\text{Gr}_I(K^\bullet \otimes_A M)
$$
Since $K^n$ is finite free we have
$$
\text{Gr}_I(K^\bullet \otimes_A M) =
\text{Gr}_I(K^\bullet) \otimes_{\text{Gr}_I(A)} \text{Gr}_I(M)
$$
Note that $\text{Gr}_I(K^\bullet)$ is the Koszul
complex over $\text{Gr}_I(A)$ on the elements
$\overline{f}_1, \ldots, \overline{f}_r \in I/I^2$.
A simple calculation (omitted)
shows that the differential $d_0$ on $E_0$
agrees with the differential coming from the Koszul complex.
Since $\text{Gr}_I(M)$ is a finite $\text{Gr}_I(A)$-module
and since $\text{Gr}_I(A)$ is Noetherian (as a quotient
of $A/I[x_1, \ldots, x_r]$ with $x_i \mapsto \overline{f}_i$), the
cohomology module $E_1 = \bigoplus E_1^{p, q}$
is a finite $\text{Gr}_I(A)$-module. However, as above
$E_1$ is annihilated by $\overline{f}_1, \ldots, \overline{f}_r$.
We conclude $E_1$ has finite length.
In particular we find that $\text{Gr}^p_F(K^\bullet \otimes M)$ is
acyclic for $p \gg 0$.

\medskip\noindent
Next, we check that the spectral sequence above converges
using Homology, Lemma \ref{homology-lemma-filtered-complex-ss-converges}.
The required equalities follow easily from the Artin-Rees lemma
in the form stated in Algebra, Lemma \ref{algebra-lemma-map-AR}.
Thus we see that
\begin{align*}
\sum (-1)^i\text{length}_A(H^i(K^\bullet \otimes_A M))
& =
\sum (-1)^{p + q} \text{length}_A(E_\infty^{p, q}) \\
& =
\sum (-1)^{p + q} \text{length}_A(E_1^{p, q})
\end{align*}
because as we've seen above the length of $E_1$ is finite
(of course this uses additivity of lengths). Pick $t$ so
large that $\text{Gr}^p_F(K^\bullet \otimes M)$
is acyclic for $p \geq t$ (see above). Using
additivity again we see that
$$
\sum (-1)^{p + q} \text{length}_A(E_1^{p, q}) =
\sum\nolimits_n \sum\nolimits_{p \leq t}
(-1)^n \text{length}_A(\text{gr}^p(K^n \otimes_A M))
$$
This is equal to
$$
\sum\nolimits_{n = -r, \ldots, 0} (-1)^n{r \choose |n|} \chi_{I, M}(t + n)
$$
by our choice of filtration above and the definition of $\chi_{I, M}$ in
Algebra, Section \ref{algebra-section-Noetherian-local}.
The lemma follows from Lemma \ref{lemma-leading-coefficient}
and the definition of $e_I(M, r)$.
\end{proof}

\bibliographystyle{plain}
\bibliography{my}
\end{document}
